\section{Supplementary Reproducibility Material}\label{sec:supplementary}
\subsection{Artifact-based reproduction}
The report is reproducible from the saved numerical tables and existing figures without fitting models. From the repository root, \code{python report/build_tables.py} regenerates the report's table fragments and source map. \code{python report/check_evidence.py} independently checks stored numerical evidence, including all h24 prediction columns, without executing notebooks. The complete-report validator checks inputs, labels, citations and source hashes. Compilation commands and the precise scope of each helper are documented in \code{report/README.md}.

The evidence extracts retain PDF page boundaries and notebook cell indices. \code{source_manifest.json} records hashes of the audited project artifacts, \code{table_sources.json} maps generated tables to CSV/JSON inputs, and \code{research_audit.md} maps important narrative claims. The report-generation helpers write within \code{report/}; they do not regenerate project experiments. The notebook runner is intentionally excluded from this compilation workflow because normal execution would fit models and overwrite saved notebook/results outputs.

\subsection{Recorded environment and experiment execution}
\Cref{tab:dependencies} reproduces the dependency pins recorded by the repository. Its experiment environment is documented as Python 3.12 with seed 42. The report inspection environment is distinct and was not used to refit the results. Package pins are provenance statements; this report does not certify that a fresh environment was installed or every artifact was reproduced end to end.
\begin{table}[H]\centering\small
\caption{Recorded experiment dependencies from requirements.txt; Python 3.12 is specified by the repository. These are recorded pins, not a newly reproduced environment.}\label{tab:dependencies}
\setlength{\tabcolsep}{4pt}
\begin{tabular}{ll}
\toprule
Package & Pinned version \\
\midrule
numpy & 2.5.3 \\
pandas & 3.0.6 \\
scipy & 1.18.1 \\
scikit-learn & 1.9.1 \\
xgboost & 3.4.1 \\
statsmodels & 0.15.0 \\
shap & 0.52.0 \\
matplotlib & 3.11.2 \\
seaborn & 0.13.2 \\
jupyter & 1.1.1 \\
ipykernel & 7.4.0 \\
nbformat & 5.11.1 \\
nbclient & 0.11.0 \\
\bottomrule\end{tabular}
\end{table}


The documented notebook order is EDA, features, November baselines, temporal CV, tuning, ablation/horizon, final test analysis and leakage diagnostics. \code{TETOUAN_N_JOBS} controls search parallelism. Smoke mode changes periods and model/search budgets and writes separate outputs under \code{results/smoke/}; its scores are code-check outputs, not the reported experiments. No full or smoke experiments were executed to prepare this report. Hardware, exact execution history, model binaries and h1 individual predictions are not archived.

\subsection{Repository organization}
\begin{verbatim}
Machine-Learning/
  README.md                 project overview
  requirements.txt          recorded experiment dependencies
  run_all.py                notebook execution runner
  data/                     existing data artifacts
  docs/                     guideline, literature review and logs
  notebooks/                ordered notebooks 01--08
  src/                      config, data, features, evaluate,
                            models, viz, and package initializer
  results/
    tables/                 saved CSV metrics/searches and JSONs
    figures/                existing PNG assets
  report/
    main.tex                root-relative LaTeX entry point
    references.bib          traceable bibliography
    research_audit.md       evidence and discrepancy audit
    sections/               thirteen report sections
    appendices/             supplementary material
    tables/                 generated LaTeX table fragments
    evidence/               extracts, manifests and checks
\end{verbatim}

\subsection{Supplementary metric variability}
\Cref{tab:cv_variability} gives mean and sample SD for every default-CV metric. These statistics describe temporal fold variation, not independent-seed repetitions. They complement the fold-specific RMSEs in \cref{tab:cv_folds}.
\begin{table}[tbp]\centering\small
\caption{Supplementary default CV variability for all four metrics (T05); mean $\pm$ sample SD. RMSE/MAE in kW, MAPE in percent.}\label{tab:cv_variability}
\setlength{\tabcolsep}{4pt}
\begin{tabular}{lrrrr}
\toprule
Model & RMSE & MAE & MAPE & R$^2$ \\
\midrule
LR & \num{2792.3} $\pm$ \num{623.3} & \num{2101.0} $\pm$ \num{468.9} & \num{2.949} $\pm$ \num{0.538} & \num{0.965} $\pm$ \num{0.016} \\
RF & \num{3069.5} $\pm$ \num{1488.8} & \num{2168.6} $\pm$ \num{985.0} & \num{2.824} $\pm$ \num{0.820} & \num{0.958} $\pm$ \num{0.030} \\
XGB & \num{3264.8} $\pm$ \num{1815.8} & \num{2328.0} $\pm$ \num{1161.3} & \num{2.992} $\pm$ \num{0.967} & \num{0.951} $\pm$ \num{0.043} \\
Naive-24h & \num{3392.4} $\pm$ \num{421.5} & \num{2371.7} $\pm$ \num{351.6} & \num{3.285} $\pm$ \num{0.318} & \num{0.950} $\pm$ \num{0.009} \\
SVR & \num{3968.7} $\pm$ \num{2891.5} & \num{2753.4} $\pm$ \num{1890.7} & \num{3.445} $\pm$ \num{1.591} & \num{0.922} $\pm$ \num{0.093} \\
DT & \num{4048.2} $\pm$ \num{1410.4} & \num{2915.3} $\pm$ \num{996.7} & \num{3.873} $\pm$ \num{0.768} & \num{0.929} $\pm$ \num{0.032} \\
Naive-168h & \num{4892.6} $\pm$ \num{2297.1} & \num{3655.4} $\pm$ \num{1673.3} & \num{4.923} $\pm$ \num{1.888} & \num{0.889} $\pm$ \num{0.077} \\
MeanProfile & \num{12005.1} $\pm$ \num{5972.6} & \num{10845.5} $\pm$ \num{5896.5} & \num{13.830} $\pm$ \num{5.454} & \num{0.375} $\pm$ \num{0.454} \\
\bottomrule\end{tabular}
\end{table}

\FloatBarrier

\subsection{Selected parameters and search traceability}
\Cref{tab:parameters} lists both horizons' selected absolute and residual configurations. Floating values are rounded for display; full precision remains in \code{best_params.json} and \code{best_params_h1.json}. A prefix omitted from the table belongs to an estimator wrapper, not to a different parameter. \Cref{tab:selected_h24} provides the compact h24 view used in the main analysis.
\begingroup\small\setlength{\tabcolsep}{4pt}
\begin{longtable}{lp{.20\linewidth}p{.32\linewidth}l}
\caption{All selected configurations from the two parameter JSON files. Wrapper prefixes are omitted; floating parameters are rounded to six decimal places here. JSON files retain full precision.}\label{tab:parameters}\\
\toprule
Setting & Variant & Parameter & Selected value \\
\midrule\endfirsthead
\multicolumn{4}{l}{\small\itshape Continued from previous page}\\\toprule
Setting & Variant & Parameter & Selected value \\
\midrule\endhead\bottomrule\endfoot
h24 & LR & alpha & \num{0.001329} \\
h24 & LR $\Delta$ & alpha & \num{98.777003} \\
h24 & DT & min\_samples\_split & 10 \\
h24 & DT & min\_samples\_leaf & 5 \\
h24 & DT & max\_features & None \\
h24 & DT & max\_depth & 8 \\
h24 & DT $\Delta$ & min\_samples\_split & 20 \\
h24 & DT $\Delta$ & min\_samples\_leaf & 5 \\
h24 & DT $\Delta$ & max\_features & None \\
h24 & DT $\Delta$ & max\_depth & 4 \\
h24 & RF & n\_estimators & 200 \\
h24 & RF & min\_samples\_leaf & 2 \\
h24 & RF & max\_features & \num{0.800000} \\
h24 & RF & max\_depth & 30 \\
h24 & RF $\Delta$ & n\_estimators & 100 \\
h24 & RF $\Delta$ & min\_samples\_leaf & 10 \\
h24 & RF $\Delta$ & max\_features & sqrt \\
h24 & RF $\Delta$ & max\_depth & 30 \\
h24 & SVR & C & \num{26.926469} \\
h24 & SVR & epsilon & \num{0.010000} \\
h24 & SVR & gamma & \num{0.006339} \\
h24 & SVR $\Delta$ & C & \num{138.262322} \\
h24 & SVR $\Delta$ & epsilon & \num{0.100000} \\
h24 & SVR $\Delta$ & gamma & \num{0.003387} \\
h24 & XGB & colsample\_bytree & \num{0.798899} \\
h24 & XGB & learning\_rate & \num{0.024629} \\
h24 & XGB & max\_depth & 3 \\
h24 & XGB & min\_child\_weight & 8 \\
h24 & XGB & n\_estimators & 1304 \\
h24 & XGB & reg\_lambda & \num{0.341635} \\
h24 & XGB & subsample & \num{0.990646} \\
h24 & XGB $\Delta$ & colsample\_bytree & \num{0.831312} \\
h24 & XGB $\Delta$ & learning\_rate & \num{0.011137} \\
h24 & XGB $\Delta$ & max\_depth & 3 \\
h24 & XGB $\Delta$ & min\_child\_weight & 10 \\
h24 & XGB $\Delta$ & n\_estimators & 774 \\
h24 & XGB $\Delta$ & reg\_lambda & \num{9.197754} \\
h24 & XGB $\Delta$ & subsample & \num{0.935573} \\
h1 & LR & alpha & \num{0.001329} \\
h1 & LR $\Delta$ & alpha & \num{0.001329} \\
h1 & DT & min\_samples\_split & 2 \\
h1 & DT & min\_samples\_leaf & 5 \\
h1 & DT & max\_features & \num{0.800000} \\
h1 & DT & max\_depth & 15 \\
h1 & DT $\Delta$ & min\_samples\_split & 20 \\
h1 & DT $\Delta$ & min\_samples\_leaf & 1 \\
h1 & DT $\Delta$ & max\_features & \num{0.500000} \\
h1 & DT $\Delta$ & max\_depth & 20 \\
h1 & RF & n\_estimators & 500 \\
h1 & RF & min\_samples\_leaf & 1 \\
h1 & RF & max\_features & \num{0.800000} \\
h1 & RF & max\_depth & 30 \\
h1 & RF $\Delta$ & n\_estimators & 300 \\
h1 & RF $\Delta$ & min\_samples\_leaf & 2 \\
h1 & RF $\Delta$ & max\_features & \num{0.500000} \\
h1 & RF $\Delta$ & max\_depth & 30 \\
h1 & SVR & C & \num{26.926469} \\
h1 & SVR & epsilon & \num{0.010000} \\
h1 & SVR & gamma & \num{0.006339} \\
h1 & SVR $\Delta$ & C & \num{44.677528} \\
h1 & SVR $\Delta$ & epsilon & \num{0.050000} \\
h1 & SVR $\Delta$ & gamma & \num{0.005019} \\
h1 & XGB & colsample\_bytree & \num{0.815877} \\
h1 & XGB & learning\_rate & \num{0.057994} \\
h1 & XGB & max\_depth & 3 \\
h1 & XGB & min\_child\_weight & 9 \\
h1 & XGB & n\_estimators & 1380 \\
h1 & XGB & reg\_lambda & \num{4.021555} \\
h1 & XGB & subsample & \num{0.629820} \\
h1 & XGB $\Delta$ & colsample\_bytree & \num{0.610168} \\
h1 & XGB $\Delta$ & learning\_rate & \num{0.013816} \\
h1 & XGB $\Delta$ & max\_depth & 9 \\
h1 & XGB $\Delta$ & min\_child\_weight & 7 \\
h1 & XGB $\Delta$ & n\_estimators & 1464 \\
h1 & XGB $\Delta$ & reg\_lambda & \num{0.425316} \\
h1 & XGB $\Delta$ & subsample & \num{0.803428} \\
\end{longtable}\endgroup


\Cref{tab:search_inventory} identifies the selected zero-based data row and mean CV RMSE for each saved search. It supports tracing a result back to individual sampled configurations. The selected search score is conditional on the same folds used to choose it; it is not an independently estimated generalization error.
\begingroup\small\setlength{\tabcolsep}{4pt}
\begin{longtable}{lrrr}
\caption{Audit of all saved searches: number of sampled configurations and selected mean CV RMSE (kW). The row index is zero-based within the CSV data, excluding its header.}\label{tab:search_inventory}\\
\toprule
Search & Candidates & Selected row & Mean RMSE \\
\midrule\endfirsthead
\multicolumn{4}{l}{\small\itshape Continued from previous page}\\\toprule
Search & Candidates & Selected row & Mean RMSE \\
\midrule\endhead\bottomrule\endfoot
LR & \num{20} & \num{10} & \num{2792.3} \\
LR\_delta & \num{20} & \num{12} & \num{2770.0} \\
DT & \num{40} & \num{17} & \num{3589.7} \\
DT\_delta & \num{40} & \num{26} & \num{2524.9} \\
RF & \num{30} & \num{25} & \num{3010.4} \\
RF\_delta & \num{30} & \num{14} & \num{2429.3} \\
SVR & \num{30} & \num{12} & \num{2410.6} \\
SVR\_delta & \num{30} & \num{10} & \num{2427.1} \\
XGB & \num{50} & \num{23} & \num{3009.9} \\
XGB\_delta & \num{50} & \num{48} & \num{2441.0} \\
h1\_LR & \num{20} & \num{10} & \num{2420.5} \\
h1\_LR\_delta & \num{20} & \num{10} & \num{2420.5} \\
h1\_DT & \num{40} & \num{34} & \num{3273.1} \\
h1\_DT\_delta & \num{40} & \num{29} & \num{2489.6} \\
h1\_RF & \num{30} & \num{21} & \num{2650.8} \\
h1\_RF\_delta & \num{30} & \num{22} & \num{2183.5} \\
h1\_SVR & \num{30} & \num{12} & \num{1866.4} \\
h1\_SVR\_delta & \num{30} & \num{18} & \num{1874.1} \\
h1\_XGB & \num{50} & \num{10} & \num{2515.1} \\
h1\_XGB\_delta & \num{50} & \num{16} & \num{2111.7} \\
\end{longtable}\endgroup


\subsection{Public-code survey}
\Cref{tab:survey} preserves the ten links and heterogeneous reported values from the repository-authored GitHub survey. No external project was rerun. Some survey descriptions simplify protocol details; the linked repositories and the local notes remain the appropriate starting points for further inspection. These entries are contextual evidence, not additional benchmark participants.
\begingroup\small\setlength{\tabcolsep}{4pt}
\begin{longtable}{p{.20\linewidth}p{.22\linewidth}p{.26\linewidth}p{.23\linewidth}}
\caption{Repository-authored survey of ten public GitHub projects (docs/kaggle\_survey.csv / T11). Metrics and protocol descriptions are survey transcriptions, not independently reproduced results; targets and metrics vary.}\label{tab:survey}\\
\toprule
Repository owner & Target & Split / validation & Survey-reported value \\
\midrule\endfirsthead
\multicolumn{4}{l}{\small\itshape Continued from previous page}\\\toprule
Repository owner & Target & Split / validation & Survey-reported value \\
\midrule\endhead\bottomrule\endfoot
\href{https://github.com/Kevinbui16/Electricity-Consumption-Prediction}{Kevinbui16} & Total (also each zone) & train\_test\_split 75/25 & 2,624 / 1,635 kW \\
\href{https://github.com/niharkaria/Tetouan_Power_Consumption_Analysis}{niharkaria} & Total & train\_test\_split 80/20 & 0.9859 / 2,026 kW \\
\href{https://github.com/nayana142/Power-Consumption-of-Tetouan-City}{nayana142} & Zone 3 & train\_test\_split 80/20 & 0.956 / 1,384 kW \\
\href{https://github.com/Merci93/Electric-Power-Consumption}{Merci93} & Zone 1, 2, 3 (separate models) & train\_test\_split 70/30 & 0.821 / 0.850 / 0.913 \\
\href{https://github.com/SoukainaBoufousse/ML_predict-the-power-consumption-in-10-minute-windows-of-Tetouan-city}{SoukainaBoufousse} & Zone 1, 2, 3 (separate models) & train\_test\_split 90/10 (Optuna validation) & 352.7 kW \\
\href{https://github.com/ThiagoLevis/Power-consumption-of-Tetouan}{ThiagoLevis} & Mean of 3 zones (min-max scaled) & train\_test\_split 70/30 & 0.19 -- 0.29 \\
\href{https://github.com/ruwindarowel/Tetouan-City-Power-Consumption}{ruwindarowel} & Total & RepeatedKFold 10x3 (CV) + hold-out & 9,907 kW \\
\href{https://github.com/atishk28/electricity-load-forecasting}{atishk28} & Total & chronological 80/20 (last 20\% as test) & 0.834 / 5,744 kW \\
\href{https://github.com/daniahfarooq/electricity-demand-forecasting}{daniahfarooq} & Zone 1, 2, 3 (per-zone models) & chronological 80/20 & 1,039 kW / 2.54\% / 0.973 \\
\href{https://github.com/AtharvaPandey-coder/Power-Consumption}{AtharvaPandey-coder} & Zone 1 & chronological 80/20 & 0.971 \\
\end{longtable}\endgroup


\subsection{Additional figures}\label{sec:supp_figures}
\Cref{fig:zone_event,fig:wind,fig:holidays} retain the quality and calendar diagnostics discussed in \cref{sec:data}. The zone event does not establish a switching mechanism, the wind distribution does not prove sensor malfunction, and holiday profiles do not control seasonal confounding. All observations are retained in the implemented preprocessing.
\begin{figure}[tbp]\centering
\includegraphics[width=\linewidth]{DQ01_zone_transfer_20170420.png}
\caption{Zone and total-load behavior around 20 April (DQ01). Redistribution between networks is a plausible interpretation; no independent operational log confirms it.}\label{fig:zone_event}
\end{figure}
\begin{figure}[tbp]\centering
\includegraphics[width=.88\linewidth]{DQ02_windspeed_bimodal.png}
\caption{Whole-year wind-speed diagnostic (DQ02). The bimodal distribution motivates caution about sensor semantics; the pipeline retains the observed values. This descriptive graphic includes December.}\label{fig:wind}
\end{figure}
\begin{figure}[tbp]\centering
\includegraphics[width=\linewidth]{F08_weekday_holiday.png}
\caption{Development-period weekday and holiday load profiles (F08). Group differences motivate calendar predictors but are neither causal effects nor independently validated event forecasts.}\label{fig:holidays}
\end{figure}
\FloatBarrier

\Cref{fig:test_scatter} supplements the selected-SVR test analysis: most observations are close to the equality line, but local departures remain. \Cref{fig:shap_dependence} shows conditional associations within the residual RF. Its lag-24 contributions explain the learned correction and exclude the explicit lag base; its temperature pattern is not a causal demand-response estimate.
\begin{figure}[tbp]\centering
\includegraphics[width=.70\linewidth]{F17c_scatter.png}
\caption{Selected h24 SVR forecasts against actual December loads (F17c). The equality reference aids assessment of large local misses; proximity alone does not establish calibrated uncertainty or temporal independence.}\label{fig:test_scatter}
\end{figure}
\begin{figure}[tbp]\centering
\includegraphics[width=\linewidth]{F19c_shap_dependence.png}
\caption{Residual-RF Tree SHAP dependence plots for temperature and lag-24 on December inputs (F19c). Values are contributions to the residual model output in kW, not to the separately added baseline.}\label{fig:shap_dependence}
\end{figure}
\FloatBarrier

\subsection{Unavailable artifacts and unresolved metadata}
The planned F09 pipeline graphic is absent; no replacement experiment output is invented. F14 is reviewed but omitted because its jointly varying parameters do not provide controlled single-parameter sensitivity. F20 is omitted to avoid visually ranking incompatible literature tasks. The related-work table and survey preserve the underlying contextual evidence instead.

Outstanding metadata include student/team names and IDs, instructor and submission date, wind/radiation units, timestamp timezone, measurement latency and detailed hardware configuration. The local row count differs from the dataset landing-page count. The literature review leaves certain paper targets, horizons and optimizer descriptions unresolved; these qualifications are retained in \cref{sec:related}. None prevents compilation, but they constrain the scientific interpretation and should be resolved before stronger claims or operational use.
